\documentclass[10pt,a4paper]{amsart}
\usepackage[margin=19mm]{geometry}
\usepackage{amsmath,amssymb,mathtools}
\usepackage[scr=rsfs,cal=euler]{mathalfa}
\usepackage{xfrac}
\usepackage[unicode,hidelinks,bookmarks=false]{hyperref}
\newcommand{\ie}{i.e.\ }
\newcommand{\cf}{cf.\ }
\newcommand{\resp}{respectively\ }
\newcommand{\Subsection}{\subsection}
\providecommand{\nobreakdash}{\nobreak\hbox{-}}
\setlength{\emergencystretch}{2em}
\multlinegap0pt
\usepackage{bm}
\usepackage{bbm}
\usepackage{dsfont}
\usepackage{xspace}
\usepackage{cancel}
\usepackage{mathtools}
\usepackage{amsthm}
\makeatletter
\@ifundefined{theorem}{\newtheorem{theorem}{Theorem}}{}
\makeatother
\title[Sharp Bounds on Diminished Sombor Index]{Sharp Bounds on Diminished Sombor Index}
\author{Meysam Taheri-Dehkordi, Amir Hossein Nokhodkar, Gholam Hossein Fath-Tabar}
\date{Source: arXiv:2607.20562v1 (CC BY 4.0)}
\begin{document}
\maketitle

\noindent\emph{Source and licence.} Sharp Bounds on Diminished Sombor Index. Version arXiv:2607.20562v1 (CC BY 4.0), distributed under the Creative Commons Attribution 4.0 International licence, as recorded in the arXiv licence metadata for this exact version and shown on the abstract page https://arxiv.org/abs/2607.20562v1. Full rights notice: licenses/arxiv-2607.20562-CC-BY-4.0.txt.\par\smallskip

\noindent\textit{Editorial scope.} Counted unit: the Theorem whose source label is \texttt{alt} in the article (source0.tex lines 931--937, source label \texttt{alt}), delivered together with the article's own complete original proof of it (source0.tex lines 938--1002, one \texttt{proof} environment of 65 lines). No mathematics is written, repaired or re-derived: statement and proof are the article's own text line for line. Required context retained uncounted: none. Every definition, display and label of those lines is retained unchanged. Omitted from this package and not redistributed: 1--930, 1003--1492. Nothing inside a retained excerpt is omitted or altered: every retained body is delivered exactly as the article's lines are written, so no omission receipt of any kind accompanies this package. Citations: the retained excerpts carry no citation command at all, so no bibliography entry of the article is transcribed and no reference is redistributed. Reference adaptation: none was needed and none was made; every reference inside the delivered unit resolves inside it, so no unresolved reference or citation is produced. Printed slips kept literal: no printed word, symbol, hypothesis or number of the counted statement or of its proof was altered. Layout and class: the article's own class is \texttt{article} with packages including amsmath, amssymb, amsthm, url, hyperref, authblk; the wrapper is the lane's pinned \texttt{amsart} editorial wrapper at 10pt on A4, which restates the theorem-like environments the retained text needs and adds the article's own macro definitions that the retained text needs (none), with extra wrapper packages (bm, bbm, dsfont, xspace, cancel, mathtools). The delivered numbering is the unit's own. Assets: no retained span contains a figure environment, a graphics-inclusion command, a PDF-page inclusion, a table or a TikZ picture; no image file is copied into this package, the package declares no asset dependency and no image file is redistributed with it. Licence: version arXiv:2607.20562v1 of this article is distributed under the Creative Commons Attribution 4.0 International licence, as recorded in the arXiv licence metadata for this exact version and shown on the abstract page https://arxiv.org/abs/2607.20562v1. A full rights notice is delivered at licenses/arxiv-2607.20562-CC-BY-4.0.txt. Characters: every retained excerpt is pure ASCII, so the delivered engine, XeLaTeX with its default Latin Modern text font, reports no missing character.
\begin{theorem}\label{alt}
Let $G$ be a  graph with $m>0$ edges. Then
\[
DSO(G) \le \sqrt{\frac{m^{2}}{2}+\frac{m}{8}\bigl(F(G)-2M_{2}(G)\bigr)}.
\]
Equality holds if and only if $G$ is regular.
\end{theorem}

\begin{proof}[Proof of the alternative bound]
By the Cauchy--Schwarz inequality,
\[
DSO(G)^{2} \le m\sum_{uv\in E(G)} \frac{d_u^{2}+d_v^{2}}{(d_u+d_v)^{2}}.
\]
For each edge,
\[
\frac{d_u^{2}+d_v^{2}}{(d_u+d_v)^{2}}
= \frac12+\frac12\cdot\frac{(d_u-d_v)^{2}}{(d_u+d_v)^{2}}.
\]
Summation over all edges yields
\[
\sum_{uv\in E(G)}\frac{d_u^{2}+d_v^{2}}{(d_u+d_v)^{2}}
= \frac{m}{2}+\frac12\sum_{uv\in E(G)}\frac{(d_u-d_v)^{2}}{(d_u+d_v)^{2}}.
\]
Because $d_u,d_v\ge 1$, we have $d_u+d_v\ge 2$ and thus $(d_u+d_v)^{2}\ge 4$.
Consequently,
\[
\frac{(d_u-d_v)^{2}}{(d_u+d_v)^{2}} \le \frac{(d_u-d_v)^{2}}{4}
\]
for every edge. Summing gives
\[
\sum_{uv\in E(G)}\frac{(d_u-d_v)^{2}}{(d_u+d_v)^{2}}
\le \frac14\sum_{uv\in E(G)}(d_u-d_v)^{2}
= \frac{\sigma(G)}{4}.
\]
Therefore,
\[
\sum_{uv\in E(G)}\frac{d_u^{2}+d_v^{2}}{(d_u+d_v)^{2}}
\le \frac{m}{2}+\frac{\sigma(G)}{8}.
\]
Plugging this into the Cauchy--Schwarz bound,
\[
DSO(G)^{2} \le m\Bigl(\frac{m}{2}+\frac{\sigma(G)}{8}\Bigr)
= \frac{m^{2}}{2}+\frac{m}{8}\,\sigma(G).
\]
Now $\sigma(G)=\sum_{uv\in E(G)}(d_u^{2}+d_v^{2}-2d_ud_v)=F(G)-2M_{2}(G)$, so we obtain the claimed inequality.

\medskip
\noindent\textit{Equality analysis.}
Assume equality holds in the final bound.
Then equality must hold in both the Cauchy--Schwarz inequality and the estimate
\[
\frac{(d_u-d_v)^{2}}{(d_u+d_v)^{2}} \le \frac{(d_u-d_v)^{2}}{4} \qquad\text{for every edge}.
\]
Equality in
\[
\frac{(d_u-d_v)^2}{(d_u+d_v)^2}
\le
\frac{(d_u-d_v)^2}{4},
\]
holds if and only if either $d_u=d_v$ (when both sides are zero) or
$(d_u+d_v)^2=4$.
Since $d_u,d_v\ge1$, the latter also implies $d_u=d_v=1$.
Hence equality is equivalent to $d_u=d_v$.
Hence every edge joins vertices of equal degree. Since $G$ is connected,
all vertices have the same degree, and therefore $G$ is regular. For any such edge with $d_u=d_v=r$, we have
\[
\frac{\sqrt{d_u^{2}+d_v^{2}}}{d_u+d_v} = \frac{\sqrt{2r^{2}}}{2r} = \frac{1}{\sqrt{2}},
\]
a constant independent of the edge. Thus the terms in the Cauchy--Schwarz sum are all equal, and equality in the Cauchy--Schwarz inequality is automatically satisfied.
Consequently, equality holds precisely for graphs whose components are regular.
For such graphs $F(G)-2M_{2}(G)=0$, and the bound reduces to $m/\sqrt{2}$, which equals $DSO(G)$.
This completes the characterization.
\end{proof}

\end{document}
